\chapter{Proof of Theorem}\label{chap:proof_of_main_theorem}

With the elimination of the $(2,1)$-cable in the previous section, we have now assembled all the necessary components to definitively prove the primary result of this dissertation. 

\begin{proof}[Proof of Theorem \ref{thm:main_detection}]
    Let $K \subset S^3$ be a knot such that $\widehat{HFK}(K) \cong \widehat{HFK}(T_{2,3} \# T_{2,3})$ as bigraded vector spaces. We must show that $K$ is isotopic to $T_{2,3} \# T_{2,3}$.
    
    First, the Euler characteristic of $\widehat{HFK}(K)$ determines the Alexander polynomial. Thus, $\Delta_K(t) = (t^2-t+1)^2$. Furthermore, by the detection results of Ni \sidecite{Ni07} and Juh\'{a}sz \sidecite{Juh08}, the top Alexander grading of $\widehat{HFK}$ forces $K$ to be a fibered knot of genus exactly $g = 2$. Let $h \colon \Sigma_{2,1} \to \Sigma_{2,1}$ be its monodromy.
    
    By the Nielsen-Thurston classification, $h$ must be freely isotopic to a periodic, pseudo-Anosov, or reducible mapping class. This partitions $K$ into three mutually exclusive geometric categories: Torus knots, Hyperbolic knots, and Satellite knots. We systematically eliminate the alternatives:
    
    \textbf{Case 1: Torus Knots.} If $h$ is periodic, $K$ is a torus knot $T_{p,q}$. As established, torus knots are L-space knots, which require $\operatorname{rank} \widehat{HFK}(K, a) \le 1$ in all Alexander gradings. However, the next-to-top Alexander grading of $\widehat{HFK}(T_{2,3} \# T_{2,3})$ has rank strictly greater than 1, establishing an immediate contradiction. $K$ cannot be a torus knot.
    
    \textbf{Case 2: Hyperbolic Knots.} If $h$ is pseudo-Anosov, $K$ is a hyperbolic knot. Furthermore, the rank configuration of $\widehat{HFK}$ restricts the fixed-point data of the monodromy, forcing its canonical lift to $\Sigma_{2,2}$ to be fixed-point-free. In Chapters \ref{chap:stratoriented} and \ref{chap:stratatraintrackanalysis}, we exhaustively constructed the state space of such mapping classes using the \texttt{autofolder} algorithm. By systematically applying the Trace Lemma, the Absorption Obstruction, and homological filtering (checking the Alexander polynomial and bounding manifold constraints), we proved that no such pseudo-Anosov map exists. Thus, $K$ cannot be hyperbolic.
    
    \textbf{Case 3: Satellite Knots.} Since $h$ is neither periodic nor pseudo-Anosov, it must be reducible, meaning $K$ is a satellite knot. In Section \ref{chap:satellite}, we utilized Schubert's formula and the classification of Thurston reducing systems to prove that the only allowable satellite configurations with winding number $|w| = 1$ and matching Alexander polynomials are the $(2,1)$-cable of the trefoil and the connect-sum of two trefoils. 
    
    In Section \ref{sec:2_1_cable_proof}, we demonstrated that the $(2,1)$-cable exhibits the incorrect Alexander polynomial ($\Delta = t^2-1+t^{-2}$), isolating the composite knot as the sole remaining topological possibility. Thus, $K$ must be the connect-sum of two trefoils.
    
    \textbf{Chirality Resolution.} It remains to show that $K$ is the homochiral sum of right-handed trefoils, $T_{2,3} \# T_{2,3}$. The square knot ($T_{2,3} \# \overline{T}_{2,3}$) and the left-handed granny knot ($\overline{T}_{2,3} \# \overline{T}_{2,3}$) share the same total rank and Alexander polynomial. However, the knot Floer homology bigradings (specifically the $\tau$ invariant and the $\delta$-gradings) are asymmetric with respect to mirroring. The bigraded vector space isomorphism $\widehat{HFK}(K) \cong \widehat{HFK}(T_{2,3} \# T_{2,3})$ strictly preserves this grading shift, uniquely identifying the right-handed homochiral sum.
    
    Consequently, $K$ is isotopic to $T_{2,3} \# T_{2,3}$, completing the proof.
\end{proof}